% Part: normal-modal-logic
% Chapter: axioms-systems
% Section: soundness

\documentclass[../../../include/open-logic-section]{subfiles}

\begin{document}

\olfileid{nml}{prf}{snd}

\olsection{निर्दोषता}

ज्या !!^a{derivation} प्रणालीमध्ये !!{derive}
करता येणारे प्रत्येक सूत्र वैध असते, तिला निर्दोष
म्हणतात. मोडल प्रणाली विचारात घेताना, म्हणजे
\Ax{K} बरोबरच काही !!{formula}s $!A_1$,
\dots,~$!A_n$ यांचे दृष्टांत वापरता येणाऱ्या
!!{derivation}s विचारात घेताना, $!A_1$, \dots,
$!A_n$ सत्य असलेल्या प्रत्येक प्रतिमानात
प्रत्येक !!{derivable} सूत्र सत्य असावे, अशी अपेक्षा आहे.

\begin{thm}[निर्दोषता प्रमेय]\ollabel{thm:soundness}
  $!A_1$, \dots, $!A_n$ यांचा प्रत्येक दृष्टांत
  अनुक्रमे प्रतिमानांच्या $\mClass{C}_1$, \dots,
  $\mClass{C}_n$ या वर्गांमध्ये वैध असेल,
  तर $\Log{K}!A_1\dots !A_n
  \Proves !B$ यावरून $!B$ हे प्रतिमानांच्या
  $\mClass{C}_1 \cap \dots \cap \mClass{C}_n$
  या वर्गात वैध ठरते.
\end{thm}

\begin{proof}
  सिद्धतांच्या लांबीवर विगमन करू. संक्षेपासाठी
  $\mClass{C} =
  \mClass{C}_1 \cap \dots \cap \mClass{C}_n$ असे ठेवा.
  \begin{enumerate}
  \item विगमनाचा पाया: $!B$ ची सिद्धता लांबी~$1$
    ची असेल, तर ते सर्वतः-सत्य दृष्टांत,
    \Ax{K} चा,\iftag{prvDiamond}{ किंवा~\Dual{} चा,}{}
    अथवा $!A_1$, \dots,~$!A_n$ यांपैकी
    एखाद्याचा दृष्टांत असतो. पहिल्या प्रकरणात
    \olref[syn][tau]{prop:valid-taut} वरून $!B$
    हा $\mClass{C}$ मध्ये वैध आहे, कारण सर्वतः-सत्य
    दृष्टांत प्रतिमानांच्या \emph{कोणत्याही} वर्गात
    वैध असतात. दुसऱ्या प्रकरणात
    \olref[syn][sch]{prop:Kvalid}\iftag{prvDiamond}{ आणि
      \olref[syn][sch]{prop:Dual-valid}}{} वरून तेच
    मिळते. शेवटी तिसऱ्या प्रकरणात $!B$ हा
    $\mClass{C}_i$ मध्ये वैध आहे आणि
    $\mClass{C} \subseteq \mClass{C}_i$ आहे;
    म्हणून \olref[syn][val]{prop:subset-class}
    वरून $!B$ हा $\mClass{C}$ मध्येही वैध आहे.
  \item विगमन पायरी: $!B$ ची सिद्धता लांबी $k>1$
    ची आहे असे समजा. $!B$ हा सर्वतः-सत्य दृष्टांत,
    मागील पायरीत नमूद केलेल्या मोडल स्वयंसिद्धकांपैकी
    एखाद्याचा दृष्टांत, किंवा $!A_1$, \dots,
    $!A_n$ यांपैकी एखाद्याचा दृष्टांत असेल,
    तर मागील पायरीप्रमाणे युक्तिवाद लागू होतो.
    म्हणून आधीच्या
    $!C \lif !B$ आणि $!C$ या !!{formula}s वर
    \MP{} लावून $!B$ मिळाला असे समजा. मग
    $!C \lif !B$ आणि $!C$ यांच्या सिद्धतांची
    लांबी $<k$ आहे; विगमन गृहीतकावरून ती
    $\mClass{C}$ मध्ये वैध आहेत.
    \olref[syn][sch]{prop:soundMP} वरून $!B$
    हाही $\mClass{C}$ मध्ये वैध आहे. शेवटी
    $!C$ वर \Nec{} लावून $!B$ मिळाला असे समजा
    (म्हणजे $!B = \Box!C$). विगमन गृहीतकावरून
    $!C$ हा $\mClass{C}$ मध्ये वैध आहे आणि
    \olref[syn][val]{prop:Nec-rule} वरून $!B$ ही
    वैध आहे.
  \end{enumerate}
\end{proof}

\end{document}
